\section{Background}\label{sec:background}

This chapter defines the LFB data window, the visual-analytics terms, and the spatial and temporal units used in the analysis.

\subsection{The London Fire Brigade and the 2024--2026 Open Data Window}\label{sec:background-lfb}

The London Fire Brigade (LFB) is the statutory fire and rescue authority for Greater London. The 2024--2026 incident slice covers 33 London local-authority areas (32 boroughs plus the City of London Corporation), 293{,}646 incidents, and 102 distinct \path{IncidentStationGround} labels. That last number counts assignment areas in the data; it is not used as a current operating-station count. LFB states its first-appliance target as a \emph{pan-London average} of six minutes~\cite{lfb_attendance}. This report separately defines an author-derived incident-level diagnostic over recorded attendance times: for any subset $A$ of the canonical slice,
\begin{equation}
E(A) \;=\; \frac{\bigl|\{\, i \in A \cap R \;:\; t_i > 360\ \mathrm{s} \,\}\bigr|}{\bigl| A \cap R \bigr|},
\qquad
R = \{\, i \;:\; t_i\ \mathrm{recorded} \,\},
\label{eq:exceedance-share}
\end{equation}
where $t_i$ is the first-pump attendance time of incident $i$ in seconds. $E(A)$ reuses the six-minute value as a descriptive threshold; it is not an official per-incident failure rate or service-level agreement. Proportions are accompanied by two-sided 95\% Wilson intervals, and dashboard cells with fewer than 30 recorded times carry a descriptive small-sample flag. The interval measures finite-sample uncertainty only: it does not correct missing response times, coordinate suppression, or measurement error. The LFB Incident Records dataset~\cite{lfb_incidents}, published continuously on the London Datastore under OGL v2, releases one row per incident with 39 fields covering call, dispatch, attendance times, incident classification, and partial location information. This project uses 2024-01 to 27 February 2026; February 2026 is a partial month.

Three features of the dataset shape every downstream design choice in this report. First, incidents are classified into \emph{Fire} (12.9\,\%), \emph{False Alarm} (46.5\,\%), and \emph{Special Service} (40.6\,\%, including medical co-responses, road traffic collisions, persons-trapped, and flooding), so comparisons must retain incident type. Second, the dataset records resources used (\texttt{NumPumpsAttending}, \texttt{NumStationsWithPumpsAttending}) but not the dispatch rule: per-station capacity, ride-time policies, and call-prioritisation logic are not released. Third, an incident row is not a mobilisation timeline. These constraints bound the analytical questions the dashboard can support.

\subsection{Visual Analytics: Operating Definitions}\label{sec:background-va}

The literature review (\S\ref{sec:literature}) draws on information-visualisation, geo-visualisation, dashboard research, uncertainty visualisation, and cognitive science, whose vocabularies sometimes use the same word differently. The following operating definitions remove that ambiguity:

\begin{itemize}
    \item \textbf{Visual analytics (VA)} --- following the Thomas \& Cook 2005 framing carried by Andrienko \& Andrienko~\cite{andrienko2007} and operationalised by Munzner~\cite{munzner2009,munzner2014}, the synergetic combination of database operations, computational analysis, and interactive visualisation in service of a human analyst's reasoning task. The emphasis is on \emph{interactive reasoning support}, not on producing a static report.
    \item \textbf{Dashboard} --- following Sarikaya et al.~\cite{sarikaya2018} and Bach et al.~\cite{bach2023patterns}, a multi-view interactive visual display whose layout, interactivity, and data semantics are bounded by an explicit purpose, audience, and update cadence. ``Dashboard'' here is therefore narrower than ``visualisation'': a single map is a visualisation but not a dashboard.
    \item \textbf{Linked view / brushing} --- following Yi et al.'s \emph{Connect} interaction category~\cite{yi2007}, a multi-panel visualisation where a selection in one panel propagates highlight to the corresponding records in every other panel. The dashboard's borough $\to$ time-series $\to$ distribution propagation is the canonical instance.
    \item \textbf{Uncertainty encoding} --- following MacEachren et al.~\cite{maceachren2012}, the use of a visual variable or adjacent field to make uncertainty visible beside the estimate. The implemented dashboard shows recorded denominators, Wilson intervals, small-sample flags, partial-month status, and median-to-p95 response ranges. Missingness and coordinate/proxy limitations are reported separately because they are not sampling uncertainty.
    \item \textbf{Station-related context} --- assignment-footprint and historical public-address proximity results calculated with straight-line distance. Travel-time coverage would require road, capacity, availability, and dispatch data; the relevant location-allocation literature is reviewed in \S\ref{sec:lit-optimisation}~\cite{church1974,goldberg2004,reconfig2023,beijing2021,ocmola2018,spacesyntax2023}.
\end{itemize}

\subsection{Geospatial and Temporal Granularities}\label{sec:background-granularity}

The dashboard works across borough, ward, and point/grid geography. Greater London contains 33 local-authority areas; wards and Lower-Layer Super Output Areas (LSOAs, mean population $\approx$\,1{,}500) provide finer administrative and statistical units. The release supplies named borough/ward/LSOA fields as well as two coordinate families. The borough dashboard joins the named borough field to the GLA boundary layer~\cite{ons_boundaries}; it does not reconstruct boroughs from incident coordinates.

The two coordinate families must not be conflated. Precise per-incident coordinates (\texttt{Easting\_m}, \texttt{Northing\_m}) are valid for about 37.9\% of rows; disclosure control accounts for much of the remainder, and the missingness may be systematic. Rounded coordinates (\texttt{Easting\_rounded}, \texttt{Northing\_rounded}) are nearly complete, but rounding introduces displacement and does not guarantee unbiased fine-scale spatial estimates. Precise points therefore support subset-only density diagnostics, while rounded points can support coarser grid checks with an explicit proxy-error caveat. Neither coordinate family limits borough counts or response summaries, which use the published borough field directly. Table~\ref{tab:evidence-support} shows which fields enter each analysis.

\begin{table}[H]
\centering
\footnotesize
\begin{tabularx}{\textwidth}{>{\raggedright\arraybackslash}p{3.1cm}*{5}{>{\centering\arraybackslash}X}}
\toprule
Analysis & Borough / rounded fields & Precise coordinates & Recorded first-pump times & Matched mobilisation partition & Station-address points \\
\midrule
Incident demand ranking and monthly trends & \textcolor{teal!60!black}{\textbf{Main input}} & -- & -- & -- & -- \\
\addlinespace
Response-performance shares and distributions (Eq.~\ref{eq:exceedance-share}) & \textcolor{teal!60!black}{\textbf{Main input}} & -- & \textcolor{teal!60!black}{\textbf{Main input}} & \textcolor{blue!60!black}{\textbf{Context}} & -- \\
\addlinespace
Street-level density diagnostics & -- & \textcolor{orange!80!black}{\textbf{Subset}} & -- & -- & -- \\
\addlinespace
Appliance-level deployment patterns & -- & -- & -- & \textcolor{teal!60!black}{\textbf{Main input}} & -- \\
\addlinespace
Station-address proximity context & \textcolor{violet!70!black}{\textbf{Context}} & -- & -- & -- & \textcolor{violet!70!black}{\textbf{Context}} \\
\midrule
Travel-time and dispatch analysis & \multicolumn{5}{c}{\textcolor{red!70!black}{\textbf{Additional station, road, capacity, and operational data required (\S4.5)}}} \\
\bottomrule
\end{tabularx}
\caption{Fields used by each analysis. Main input marks the fields used in the calculation; Subset marks the precise-coordinate diagnostic; Context marks a related comparison; a dash marks an unused field.}
\label{tab:evidence-support}
\end{table}

Temporally the slice covers 2024-01-01 to 2026-02-27 and is entirely post-COVID-19. It contains 25 complete months and one partial month, February 2026, which the dashboard marks explicitly. Sharma et al.~\cite{predictive2024} show that pandemic-era model error can vary by incident type, so comparisons with pre-pandemic studies such as Taylor's 2012--2014 window~\cite{spatial2015} should not assume a uniform shift.
